\section{Systèmes de Génération et Grands Modèles de Langue Augmentés par Extraction (RAG et RALLM)}

    \subsection{Besoin et Définition}
        \begin{verbatim}
- Les modèles profonds apprennent sur les données [ref à récupérer de Lewis+20]
- Le problème des modèles génératifs : hallucination et connaissance bornée dans 
le temps [ref à récupérer de Lewis+20]
- Il y a quelques années, on a cherché à associer un base de connaissances à ces 
modèles [Lewis+20] (= définition de RAG)
- Parmis ces approches, on retrouve le cas des modèles de langue : RALLM
        \end{verbatim}


    \subsection{L'architecture RAG}

        \begin{verbatim}
Description de l'archi de façon plus avancée
- Ces modèles peuvent être utilisés pour différents types de tâches (voir tableau 
csv sur RAGRights)
- Avec le temps, l'idée de [Lewis+20] a été adaptée et augmentée de composants 
supplémentaires
- [Fan+23] propose la structure générique ... (exploiter slides Séminaire CA)
        \end{verbatim}


    \subsection{Architectures aperçues au cours du stage}

        \begin{verbatim}
Au delà de l'innovation de l'architecture, on retrouve des innovations sur plusieurs 
caractéristiques de ces modèles : [R24/04] [R16/02]
- Multimodal [Liu+25]
- Chain of Thought [Xu+24]
- Gestion BDD dynamique [ref à trouver dans R24/05]
- BDD sparse et interprétable [Prouteau]???
        \end{verbatim}


    \subsection{Architecture et Modèles utilisées au cours du stage}

        \begin{verbatim}
Dans le cadre du stage, 
- travailler sur une adaptation du modèle RAG4HN. [Tran+25] 
- _description de l'architecture RAG4HN_ + parallèle avec l'architecture générique 
de [Fan+23]
- L'extracteur de RAG4HN est mE5 [Wang+25b] (version multilingue de E5 [Whang+25a])
        \end{verbatim}


\section{Systèmes RAG de confiance}

    \subsection{Trustworthy AI}

        \begin{verbatim}
La définition d'IA de confiance (Trustworthy AI) tend à varier en fonction des situation
et des utilisations de cette expression. Danss le contexte du Machine Learning, on désigne 
l'IA de confiance comme l'ensembles des systèmes d'IA dans lesquels il est légitime (dans 
le sens humainement entendable) d'avoir confiance. Le NIST (National Institute of Standards 
and Technology (US)) propose différents axes ou composantes de l'IA de confiance : ...
        \end{verbatim}


    \subsection{Explicabilité pour le RAG}
        \begin{verbatim}
Dans la suite, nous nous concentrons sur l'explicabilité dans le contexte du RAG. [Ni+25] 
distingue deux temps pour l'explicabilité : Extraction et génération ... (prévoir la question 
de l'association/intégration des deux)
- Métrique ??
        \end{verbatim}


    \subsection{Interprétabilité Mécaniste}
        \begin{verbatim}
[Somvanshi+25]
- Interprétabilité vs Explicabilité
- Définition d'Interprétabilité Mécaniste et approches (Manual Circuit Tracing,
*Intervention-based technique*, Representation analysis, Toy Model and synthetic tasks)
        \end{verbatim}


    \subsection{Activation Patching}

        \begin{verbatim}
- Définition de activation patching [Chen+24] [R30/04] [Séminaire CA]
- MechIR [Parry+25] 
- Perturbation
- Données produites, visualisation et interprétation des graphiques
        \end{verbatim}